% Bài: L12C1 Bài 2. Nội năng. Định luật I của nhiệt động lực học · Các cách làm biến đổi nội năng của một vật
% ID: L12C1B2-01-DS
% Mức: NB
% ID: L12C1B2-01-DS
% ID: L12C1B2-01-DS
% Mức: NB
% ID: L12C1B3-01-DS
% ID: L12C1B3-01-DS
% Mức: NB
% ID: L12C1B2-02-DS
%=== Câu 1 ===%
\dangbt{Các cách làm biến đổi nội năng của một vật}
\begin{ex}
% ID: L12C1B2-02-DS
% Mức: NB
Xét các phát biểu sau về dạng “Các cách làm biến đổi nội năng của một vật”
\choiceTF
{\True Nội năng có thể thay đổi bằng thực hiện công hoặc truyền nhiệt.}
{Truyền nhiệt chỉ xảy ra khi hai vật có cùng nhiệt độ.}
{\True Khi giải dạng “Các cách làm biến đổi nội năng của một vật”, cần xác định đúng trạng thái đầu, trạng thái cuối và quy ước dấu hoặc điều kiện áp dụng.}
{Có thể bỏ qua mọi dữ kiện về khối lượng, nhiệt độ và hiệu suất trong mọi bài toán thuộc dạng này.}
\loigiai{
\begin{itemchoice}
\item (Đúng) Nội năng có thể thay đổi bằng thực hiện công hoặc truyền nhiệt. (Vì): Nội năng có thể thay đổi bằng thực hiện công hoặc truyền nhiệt. b) S: Truyền nhiệt chỉ xảy ra khi hai vật có cùng nhiệt độ. c) Đ: Khi giải dạng “Các cách làm biến đổi nội năng của một vật”, cần xác định đúng trạng thái đầu, trạng thái cuối và quy ước dấu hoặc điều kiện áp dụng. d) S: Có thể bỏ qua mọi dữ kiện về khối lượng, nhiệt độ và hiệu suất trong mọi bài toán thuộc dạng này. Kết luận dựa trên định nghĩa, điều kiện áp dụng và quy ước trong SGK KNTT
\item (Sai) Truyền nhiệt chỉ xảy ra khi hai vật có cùng nhiệt độ.
\item (Đúng) Khi giải dạng “Các cách làm biến đổi nội năng của một vật”, cần xác định đúng trạng thái đầu, trạng thái cuối và quy ước dấu hoặc điều kiện áp dụng.
\item (Sai) Có thể bỏ qua mọi dữ kiện về khối lượng, nhiệt độ và hiệu suất trong mọi bài toán thuộc dạng này.
\end{itemchoice}
}
\end{ex}

% ID: L12C1B2-01-TL
% Mức: NB
% ID: L12C1B2-01-TL
% ID: L12C1B2-01-TL
% Mức: NB
% ID: L12C1B3-02-TL
% ID: L12C1B3-02-TL
% Mức: NB
% ID: L12C1B2-03-TL
%=== Câu 2 ===%
\begin{bt}
% ID: L12C1B2-03-TL
% Mức: NB
Trình bày cách nhận biết và phương pháp giải dạng “Các cách làm biến đổi nội năng của một vật”. Phân tích nhận định sau và sửa lại nếu sai: “Truyền nhiệt chỉ xảy ra khi hai vật có cùng nhiệt độ.”
\loigiai{
Trước hết xác định hiện tượng, trạng thái và các đại lượng đã cho. Nêu định nghĩa hoặc hệ thức phù hợp, đổi về đơn vị SI, áp dụng đúng điều kiện và kiểm tra ý nghĩa kết quả. Nhận định cần đối chiếu: Nội năng có thể thay đổi bằng thực hiện công hoặc truyền nhiệt. Vì vậy nhận định đã cho sai; phát biểu đúng là: Nội năng có thể thay đổi bằng thực hiện công hoặc truyền nhiệt.
}
\end{bt}

% ID: L12C1B2-01-TN
% Mức: NB
% ID: L12C1B2-01-TN
% ID: L12C1B2-01-TN
% Mức: NB
% ID: L12C1B3-04-TN
% ID: L12C1B3-04-TN
% Mức: NB
% ID: L12C1B2-04-TN
%=== Câu 4 ===%
\begin{ex}
% ID: L12C1B2-04-TN
% Mức: NB
Câu 1 thuộc dạng “Các cách làm biến đổi nội năng của một vật”. Phát biểu nào sau đây đúng?
\choice
{Truyền nhiệt chỉ xảy ra khi hai vật có cùng nhiệt độ.}
{Nhiệt lượng và nhiệt độ là hai đại lượng hoàn toàn giống nhau.}
{Mọi quá trình nhiệt học đều không phụ thuộc điều kiện thực hiện.}
{\True Nội năng có thể thay đổi bằng thực hiện công hoặc truyền nhiệt.}
\loigiai{
Phát biểu đúng là: Nội năng có thể thay đổi bằng thực hiện công hoặc truyền nhiệt. Các phương án còn lại trái với mô hình, định nghĩa hoặc hệ thức nhiệt học tương ứng.
}
\end{ex}

% ID: L12C1B2-02-DS
% Mức: TH
% ID: L12C1B2-02-DS
% ID: L12C1B2-02-DS
% Mức: TH
% ID: L12C1B3-05-DS
% ID: L12C1B3-05-DS
% Mức: TH
% ID: L12C1B2-05-DS
%=== Câu 5 ===%
\begin{ex}
% ID: L12C1B2-05-DS
% Mức: TH
Xét các phát biểu sau về dạng “Các cách làm biến đổi nội năng của một vật”
\choiceTF
{Truyền nhiệt chỉ xảy ra khi hai vật có cùng nhiệt độ.}
{\True Khi giải dạng “Các cách làm biến đổi nội năng của một vật”, cần xác định đúng trạng thái đầu, trạng thái cuối và quy ước dấu hoặc điều kiện áp dụng.}
{Có thể bỏ qua mọi dữ kiện về khối lượng, nhiệt độ và hiệu suất trong mọi bài toán thuộc dạng này.}
{\True Nội năng có thể thay đổi bằng thực hiện công hoặc truyền nhiệt.}
\loigiai{
\begin{itemchoice}
\item (Sai) Truyền nhiệt chỉ xảy ra khi hai vật có cùng nhiệt độ. (Vì): Truyền nhiệt chỉ xảy ra khi hai vật có cùng nhiệt độ. b) Đ: Khi giải dạng “Các cách làm biến đổi nội năng của một vật”, cần xác định đúng trạng thái đầu, trạng thái cuối và quy ước dấu hoặc điều kiện áp dụng. c) S: Có thể bỏ qua mọi dữ kiện về khối lượng, nhiệt độ và hiệu suất trong mọi bài toán thuộc dạng này. d) Đ: Nội năng có thể thay đổi bằng thực hiện công hoặc truyền nhiệt. Kết luận dựa trên định nghĩa, điều kiện áp dụng và quy ước trong SGK KNTT
\item (Đúng) Khi giải dạng “Các cách làm biến đổi nội năng của một vật”, cần xác định đúng trạng thái đầu, trạng thái cuối và quy ước dấu hoặc điều kiện áp dụng.
\item (Sai) Có thể bỏ qua mọi dữ kiện về khối lượng, nhiệt độ và hiệu suất trong mọi bài toán thuộc dạng này.
\item (Đúng) Nội năng có thể thay đổi bằng thực hiện công hoặc truyền nhiệt.
\end{itemchoice}
}
\end{ex}

% ID: L12C1B2-02-TLN
% Mức: NB
% ID: L12C1B2-02-TLN
% ID: L12C1B2-02-TLN
% Mức: NB
% ID: L12C1B3-06-TLN
% ID: L12C1B3-06-TLN
% Mức: NB
% ID: L12C1B2-06-TLN
%=== Câu 6 ===%
\begin{ex}
% ID: L12C1B2-06-TLN
% Mức: NB
Cho bốn nhận định: (1) Nội năng có thể thay đổi bằng thực hiện công hoặc truyền nhiệt. (2) Truyền nhiệt chỉ xảy ra khi hai vật có cùng nhiệt độ. (3) Cần kiểm tra đơn vị và điều kiện áp dụng khi xử lí dạng “Các cách làm biến đổi nội năng của một vật”. (4) Có thể thay tùy ý nhiệt độ Celsius cho nhiệt độ tuyệt đối trong mọi công thức. Có bao nhiêu nhận định đúng? Trả lời bằng một số nguyên.
\shortans{2}
\loigiai{
Nhận định (1) và (3) đúng; (2) và (4) sai. Vậy có 2 nhận định đúng.
}
\end{ex}

% ID: L12C1B2-02-TN
% Mức: NB
% ID: L12C1B2-02-TN
% ID: L12C1B2-02-TN
% Mức: NB
% ID: L12C1B3-07-TN
% ID: L12C1B3-07-TN
% Mức: NB
% ID: L12C1B2-07-TN
%=== Câu 7 ===%
\begin{ex}
% ID: L12C1B2-07-TN
% Mức: NB
Câu 5 thuộc dạng “Các cách làm biến đổi nội năng của một vật”. Phát biểu nào sau đây đúng?
\choice
{Truyền nhiệt chỉ xảy ra khi hai vật có cùng nhiệt độ.}
{Nhiệt lượng và nhiệt độ là hai đại lượng hoàn toàn giống nhau.}
{Mọi quá trình nhiệt học đều không phụ thuộc điều kiện thực hiện.}
{\True Nội năng có thể thay đổi bằng thực hiện công hoặc truyền nhiệt.}
\loigiai{
Phát biểu đúng là: Nội năng có thể thay đổi bằng thực hiện công hoặc truyền nhiệt. Các phương án còn lại trái với mô hình, định nghĩa hoặc hệ thức nhiệt học tương ứng.
}
\end{ex}

% ID: L12C1B2-03-DS
% Mức: TH
% ID: L12C1B2-03-DS
% ID: L12C1B2-03-DS
% Mức: TH
% ID: L12C1B3-08-DS
% ID: L12C1B3-08-DS
% Mức: TH
% ID: L12C1B2-08-DS
%=== Câu 8 ===%
\begin{ex}
% ID: L12C1B2-08-DS
% Mức: TH
Xét các phát biểu sau về dạng “Các cách làm biến đổi nội năng của một vật”
\choiceTF
{Truyền nhiệt chỉ xảy ra khi hai vật có cùng nhiệt độ.}
{\True Khi giải dạng “Các cách làm biến đổi nội năng của một vật”, cần xác định đúng trạng thái đầu, trạng thái cuối và quy ước dấu hoặc điều kiện áp dụng.}
{Có thể bỏ qua mọi dữ kiện về khối lượng, nhiệt độ và hiệu suất trong mọi bài toán thuộc dạng này.}
{\True Nội năng có thể thay đổi bằng thực hiện công hoặc truyền nhiệt.}
\loigiai{
\begin{itemchoice}
\item (Sai) Truyền nhiệt chỉ xảy ra khi hai vật có cùng nhiệt độ. (Vì): Truyền nhiệt chỉ xảy ra khi hai vật có cùng nhiệt độ. b) Đ: Khi giải dạng “Các cách làm biến đổi nội năng của một vật”, cần xác định đúng trạng thái đầu, trạng thái cuối và quy ước dấu hoặc điều kiện áp dụng. c) S: Có thể bỏ qua mọi dữ kiện về khối lượng, nhiệt độ và hiệu suất trong mọi bài toán thuộc dạng này. d) Đ: Nội năng có thể thay đổi bằng thực hiện công hoặc truyền nhiệt. Kết luận dựa trên định nghĩa, điều kiện áp dụng và quy ước trong SGK KNTT
\item (Đúng) Khi giải dạng “Các cách làm biến đổi nội năng của một vật”, cần xác định đúng trạng thái đầu, trạng thái cuối và quy ước dấu hoặc điều kiện áp dụng.
\item (Sai) Có thể bỏ qua mọi dữ kiện về khối lượng, nhiệt độ và hiệu suất trong mọi bài toán thuộc dạng này.
\item (Đúng) Nội năng có thể thay đổi bằng thực hiện công hoặc truyền nhiệt.
\end{itemchoice}
}
\end{ex}

% ID: L12C1B2-03-TL
% Mức: TH
% ID: L12C1B2-03-TL
% ID: L12C1B2-03-TL
% Mức: TH
% ID: L12C1B3-09-TL
% ID: L12C1B3-09-TL
% Mức: TH
% ID: L12C1B2-09-TL
%=== Câu 9 ===%
\begin{bt}
% ID: L12C1B2-09-TL
% Mức: TH
Trình bày cách nhận biết và phương pháp giải dạng “Các cách làm biến đổi nội năng của một vật”. Phân tích nhận định sau và sửa lại nếu sai: “Nội năng có thể thay đổi bằng thực hiện công hoặc truyền nhiệt.”
\loigiai{
Trước hết xác định hiện tượng, trạng thái và các đại lượng đã cho. Nêu định nghĩa hoặc hệ thức phù hợp, đổi về đơn vị SI, áp dụng đúng điều kiện và kiểm tra ý nghĩa kết quả. Nhận định cần đối chiếu: Nội năng có thể thay đổi bằng thực hiện công hoặc truyền nhiệt. Vì vậy nhận định đã cho đúng.
}
\end{bt}

% ID: L12C1B2-03-TN
% Mức: NB
% ID: L12C1B2-03-TN
% ID: L12C1B2-03-TN
% Mức: NB
% ID: L12C1B3-10-TN
% ID: L12C1B3-10-TN
% Mức: NB
% ID: L12C1B2-10-TN
%=== Câu 10 ===%
\begin{ex}
% ID: L12C1B2-10-TN
% Mức: NB
Câu 9 thuộc dạng “Các cách làm biến đổi nội năng của một vật”. Phát biểu nào sau đây đúng?
\choice
{Truyền nhiệt chỉ xảy ra khi hai vật có cùng nhiệt độ.}
{Nhiệt lượng và nhiệt độ là hai đại lượng hoàn toàn giống nhau.}
{Mọi quá trình nhiệt học đều không phụ thuộc điều kiện thực hiện.}
{\True Nội năng có thể thay đổi bằng thực hiện công hoặc truyền nhiệt.}
\loigiai{
Phát biểu đúng là: Nội năng có thể thay đổi bằng thực hiện công hoặc truyền nhiệt. Các phương án còn lại trái với mô hình, định nghĩa hoặc hệ thức nhiệt học tương ứng.
}
\end{ex}

% ID: L12C1B2-04-DS
% Mức: VD
% ID: L12C1B2-04-DS
% ID: L12C1B2-04-DS
% Mức: VD
% ID: L12C1B3-16-DS
% ID: L12C1B3-16-DS
% Mức: VD
% ID: L12C1B2-12-DS
%=== Câu 11 ===%
\begin{ex}
% ID: L12C1B2-12-DS
% Mức: VD
Xét các phát biểu sau về dạng “Các cách làm biến đổi nội năng của một vật” 
\choiceTF
{\True Khi giải dạng “Các cách làm biến đổi nội năng của một vật”, cần xác định đúng trạng thái đầu, trạng thái cuối và quy ước dấu hoặc điều kiện áp dụng.}
{Có thể bỏ qua mọi dữ kiện về khối lượng, nhiệt độ và hiệu suất trong mọi bài toán thuộc dạng này.}
{\True Nội năng có thể thay đổi bằng thực hiện công hoặc truyền nhiệt.}
{Truyền nhiệt chỉ xảy ra khi hai vật có cùng nhiệt độ.}
\loigiai{
\begin{itemchoice}
\item (Đúng) Khi giải dạng “Các cách làm biến đổi nội năng của một vật”, cần xác định đúng trạng thái đầu, trạng thái cuối và quy ước dấu hoặc điều kiện áp dụng. (Vì): Khi giải dạng “Các cách làm biến đổi nội năng của một vật”, cần xác định đúng trạng thái đầu, trạng thái cuối và quy ước dấu hoặc điều kiện áp dụng. b) S: Có thể bỏ qua mọi dữ kiện về khối lượng, nhiệt độ và hiệu suất trong mọi bài toán thuộc dạng này. c) Đ: Nội năng có thể thay đổi bằng thực hiện công hoặc truyền nhiệt. d) S: Truyền nhiệt chỉ xảy ra khi hai vật có cùng nhiệt độ. Kết luận dựa trên định nghĩa, điều kiện áp dụng và quy ước trong SGK KNTT
\item (Sai) Có thể bỏ qua mọi dữ kiện về khối lượng, nhiệt độ và hiệu suất trong mọi bài toán thuộc dạng này.
\item (Đúng) Nội năng có thể thay đổi bằng thực hiện công hoặc truyền nhiệt.
\item (Sai) Truyền nhiệt chỉ xảy ra khi hai vật có cùng nhiệt độ.
\end{itemchoice}
}
\end{ex}

% ID: L12C1B2-04-TN
% Mức: TH
% ID: L12C1B2-04-TN
% ID: L12C1B2-04-TN
% Mức: TH
% ID: L12C1B3-17-TN
% ID: L12C1B3-17-TN
% Mức: TH
% ID: L12C1B2-13-TN
%=== Câu 12 ===%
\begin{ex}
% ID: L12C1B2-13-TN
% Mức: TH
Câu 2 thuộc dạng “Các cách làm biến đổi nội năng của một vật”. Phát biểu nào sau đây đúng?
\choice
{Nhiệt lượng và nhiệt độ là hai đại lượng hoàn toàn giống nhau.}
{Mọi quá trình nhiệt học đều không phụ thuộc điều kiện thực hiện.}
{\True Nội năng có thể thay đổi bằng thực hiện công hoặc truyền nhiệt.}
{Truyền nhiệt chỉ xảy ra khi hai vật có cùng nhiệt độ.}
\loigiai{
Phát biểu đúng là: Nội năng có thể thay đổi bằng thực hiện công hoặc truyền nhiệt. Các phương án còn lại trái với mô hình, định nghĩa hoặc hệ thức nhiệt học tương ứng.
}
\end{ex}

% ID: L12C1B2-05-DS
% Mức: VDC
% ID: L12C1B2-05-DS
% ID: L12C1B2-05-DS
% Mức: VDC
% ID: L12C1B3-19-DS
% ID: L12C1B3-19-DS
% Mức: VDC
% ID: L12C1B2-14-DS
%=== Câu 13 ===%
\begin{ex}
% ID: L12C1B2-14-DS
% Mức: VDC
Xét các phát biểu sau về dạng “Các cách làm biến đổi nội năng của một vật” 
\choiceTF
{Có thể bỏ qua mọi dữ kiện về khối lượng, nhiệt độ và hiệu suất trong mọi bài toán thuộc dạng này.}
{\True Nội năng có thể thay đổi bằng thực hiện công hoặc truyền nhiệt.}
{Truyền nhiệt chỉ xảy ra khi hai vật có cùng nhiệt độ.}
{\True Khi giải dạng “Các cách làm biến đổi nội năng của một vật”, cần xác định đúng trạng thái đầu, trạng thái cuối và quy ước dấu hoặc điều kiện áp dụng.}
\loigiai{
\begin{itemchoice}
\item (Sai) Có thể bỏ qua mọi dữ kiện về khối lượng, nhiệt độ và hiệu suất trong mọi bài toán thuộc dạng này. (Vì): Có thể bỏ qua mọi dữ kiện về khối lượng, nhiệt độ và hiệu suất trong mọi bài toán thuộc dạng này. b) Đ: Nội năng có thể thay đổi bằng thực hiện công hoặc truyền nhiệt. c) S: Truyền nhiệt chỉ xảy ra khi hai vật có cùng nhiệt độ. d) Đ: Khi giải dạng “Các cách làm biến đổi nội năng của một vật”, cần xác định đúng trạng thái đầu, trạng thái cuối và quy ước dấu hoặc điều kiện áp dụng. Kết luận dựa trên định nghĩa, điều kiện áp dụng và quy ước trong SGK KNTT
\item (Đúng) Nội năng có thể thay đổi bằng thực hiện công hoặc truyền nhiệt.
\item (Sai) Truyền nhiệt chỉ xảy ra khi hai vật có cùng nhiệt độ.
\item (Đúng) Khi giải dạng “Các cách làm biến đổi nội năng của một vật”, cần xác định đúng trạng thái đầu, trạng thái cuối và quy ước dấu hoặc điều kiện áp dụng.
\end{itemchoice}
}
\end{ex}

% ID: L12C1B2-05-TN
% Mức: TH
% ID: L12C1B2-05-TN
% ID: L12C1B2-05-TN
% Mức: TH
% ID: L12C1B3-73-TN
% ID: L12C1B3-73-TN
% Mức: TH
% ID: L12C1B2-15-TN
%=== Câu 14 ===%
\begin{ex}
% ID: L12C1B2-15-TN
% Mức: TH
Câu 6 thuộc dạng “Các cách làm biến đổi nội năng của một vật”. Phát biểu nào sau đây đúng?
\choice
{Nhiệt lượng và nhiệt độ là hai đại lượng hoàn toàn giống nhau.}
{Mọi quá trình nhiệt học đều không phụ thuộc điều kiện thực hiện.}
{\True Nội năng có thể thay đổi bằng thực hiện công hoặc truyền nhiệt.}
{Truyền nhiệt chỉ xảy ra khi hai vật có cùng nhiệt độ.}
\loigiai{
Phát biểu đúng là: Nội năng có thể thay đổi bằng thực hiện công hoặc truyền nhiệt. Các phương án còn lại trái với mô hình, định nghĩa hoặc hệ thức nhiệt học tương ứng.
}
\end{ex}

% ID: L12C1B2-06-TN
% Mức: TH
% ID: L12C1B2-06-TN
% ID: L12C1B2-06-TN
% Mức: TH
% ID: L12C1B3-74-TN
% ID: L12C1B3-74-TN
% Mức: TH
% ID: L12C1B2-16-TN
%=== Câu 16 ===%
\begin{ex}
% ID: L12C1B2-16-TN
% Mức: TH
Câu 10 thuộc dạng “Các cách làm biến đổi nội năng của một vật”. Phát biểu nào sau đây đúng?
\choice
{Nhiệt lượng và nhiệt độ là hai đại lượng hoàn toàn giống nhau.}
{Mọi quá trình nhiệt học đều không phụ thuộc điều kiện thực hiện.}
{\True Nội năng có thể thay đổi bằng thực hiện công hoặc truyền nhiệt.}
{Truyền nhiệt chỉ xảy ra khi hai vật có cùng nhiệt độ.}
\loigiai{
Phát biểu đúng là: Nội năng có thể thay đổi bằng thực hiện công hoặc truyền nhiệt. Các phương án còn lại trái với mô hình, định nghĩa hoặc hệ thức nhiệt học tương ứng.
}
\end{ex}

% ID: L12C1B2-07-TN
% Mức: VD
% ID: L12C1B2-07-TN
% ID: L12C1B2-07-TN
% Mức: VD
% ID: L12C1B3-75-TN
% ID: L12C1B3-75-TN
% Mức: VD
% ID: L12C1B2-17-TN
%=== Câu 19 ===%
\begin{ex}
% ID: L12C1B2-17-TN
% Mức: VD
Câu 3 thuộc dạng “Các cách làm biến đổi nội năng của một vật”. Phát biểu nào sau đây đúng?
\choice
{Mọi quá trình nhiệt học đều không phụ thuộc điều kiện thực hiện.}
{\True Nội năng có thể thay đổi bằng thực hiện công hoặc truyền nhiệt.}
{Truyền nhiệt chỉ xảy ra khi hai vật có cùng nhiệt độ.}
{Nhiệt lượng và nhiệt độ là hai đại lượng hoàn toàn giống nhau.}
\loigiai{
Phát biểu đúng là: Nội năng có thể thay đổi bằng thực hiện công hoặc truyền nhiệt. Các phương án còn lại trái với mô hình, định nghĩa hoặc hệ thức nhiệt học tương ứng.
}
\end{ex}

% ID: L12C1B2-08-TN
% Mức: VD
% ID: L12C1B2-08-TN
% ID: L12C1B2-08-TN
% Mức: VD
% ID: L12C1B3-76-TN
% ID: L12C1B3-76-TN
% Mức: VD
% ID: L12C1B2-18-TN
%=== Câu 22 ===%
\begin{ex}
% ID: L12C1B2-18-TN
% Mức: VD
Câu 7 thuộc dạng “Các cách làm biến đổi nội năng của một vật”. Phát biểu nào sau đây đúng?
\choice
{Mọi quá trình nhiệt học đều không phụ thuộc điều kiện thực hiện.}
{\True Nội năng có thể thay đổi bằng thực hiện công hoặc truyền nhiệt.}
{Truyền nhiệt chỉ xảy ra khi hai vật có cùng nhiệt độ.}
{Nhiệt lượng và nhiệt độ là hai đại lượng hoàn toàn giống nhau.}
\loigiai{
Phát biểu đúng là: Nội năng có thể thay đổi bằng thực hiện công hoặc truyền nhiệt. Các phương án còn lại trái với mô hình, định nghĩa hoặc hệ thức nhiệt học tương ứng.
}
\end{ex}

% ID: L12C1B2-09-TN
% Mức: VDC
% ID: L12C1B2-09-TN
% ID: L12C1B2-09-TN
% Mức: VDC
% ID: L12C1B3-77-TN
% ID: L12C1B3-77-TN
% Mức: VDC
% ID: L12C1B2-19-TN
%=== Câu 25 ===%
\begin{ex}
% ID: L12C1B2-19-TN
% Mức: VDC
Câu 4 thuộc dạng “Các cách làm biến đổi nội năng của một vật”. Phát biểu nào sau đây đúng?
\choice
{\True Nội năng có thể thay đổi bằng thực hiện công hoặc truyền nhiệt.}
{Truyền nhiệt chỉ xảy ra khi hai vật có cùng nhiệt độ.}
{Nhiệt lượng và nhiệt độ là hai đại lượng hoàn toàn giống nhau.}
{Mọi quá trình nhiệt học đều không phụ thuộc điều kiện thực hiện.}
\loigiai{
Phát biểu đúng là: Nội năng có thể thay đổi bằng thực hiện công hoặc truyền nhiệt. Các phương án còn lại trái với mô hình, định nghĩa hoặc hệ thức nhiệt học tương ứng.
}
\end{ex}

% ID: L12C1B2-10-TN
% Mức: VDC
% ID: L12C1B2-10-TN
% ID: L12C1B2-10-TN
% Mức: VDC
% ID: L12C1B3-78-TN
% ID: L12C1B3-78-TN
% Mức: VDC
% ID: L12C1B2-20-TN
%=== Câu 27 ===%
\begin{ex}
% ID: L12C1B2-20-TN
% Mức: VDC
Câu 8 thuộc dạng “Các cách làm biến đổi nội năng của một vật”. Phát biểu nào sau đây đúng?
\choice
{\True Nội năng có thể thay đổi bằng thực hiện công hoặc truyền nhiệt.}
{Truyền nhiệt chỉ xảy ra khi hai vật có cùng nhiệt độ.}
{Nhiệt lượng và nhiệt độ là hai đại lượng hoàn toàn giống nhau.}
{Mọi quá trình nhiệt học đều không phụ thuộc điều kiện thực hiện.}
\loigiai{
Phát biểu đúng là: Nội năng có thể thay đổi bằng thực hiện công hoặc truyền nhiệt. Các phương án còn lại trái với mô hình, định nghĩa hoặc hệ thức nhiệt học tương ứng.
}
\end{ex}

% Mức: NB
% ID: L12C1B2-82-TN
% ID: L12C1B2-82-TN
% Mức: NB
% ID: L12C1B3-79-TN
% ID: L12C1B3-79-TN
% Mức: NB
% ID: L12C1B2-74-TN
%=== Câu 90 ===%
\begin{ex}
% ID: L12C1B2-74-TN
% Mức: NB
	Trường hợp nào dưới đây làm biến đổi nội năng \textbf{không} phải do thực hiện công?
	\choice 
	{Cọ xát hai vật vào nhau}
	{Nén khí trong xilanh}
	{\True Đun nóng nước bằng bếp}
	{Một viên bi bằng thép rơi xuống đất mềm}
	\loigiai{
		Có hai cách làm thay đổi nội năng của vật là thực hiện công và truyền nhiệt.
		\begin{itemize}
			\item Các trường hợp: cọ xát hai vật vào nhau, nén khí trong xilanh, một viên bi bằng thép rơi xuống đất mềm đều có sự chuyển hoá năng lượng từ cơ năng sang nội năng, nên đây là quá trình làm thay đổi nội năng bằng cách thực hiện công.
			\item Trường hợp đun nóng nước bằng bếp chỉ có sự truyền năng lượng nhiệt (không có sự chuyển hoá từ cơ năng sang nội năng) nên đây là quá trình truyền nhiệt,.
		\end{itemize}
	}
\end{ex}
